% from main.tex
%%%%%%%%%%  scan idx 550  |  folio 537  %%%%%%%%%%
{\operatorname{length}_{#1}}


\subsection*{2.4, suite}

Lorsque \(X\) est propre sur \(k\), on sait [4] que \(\Pic^0_{X/k}\) est représentable par un schéma en groupes localement de type fini \(\Pic_{X/k}\) sur \(k\). La composante neutre \(\Pic^0_{X/k}\) de \(\Pic_{X/k}\) est un groupe algébrique commutatif connexe sur \(k\).

\textbf{Définition 2.4.1.} Soit \(X\) un schéma propre sur un corps \(k\), \(L\) un \(\mathcal O_X\)-module inversible (resp. \(x\) un élément de \(\Pic(X)\)). On dit que \(L\)(resp. \(x\)) est \(\tau\)-équivalent à zéro, s'il existe un entier \(n\ne0\) tel que \(L^{\otimes n}\)(resp. \(nx\)) est algébriquement équivalent à zéro. Les éléments de \(\Pic(X)\) algébriquement équivalents (resp. \(\tau\)-équivalents) à zéro forment un sous-groupe \(\Pic^0(X)\) (resp. \(\Pic^\tau(X)\)) de \(\Pic(X)\).

\textbf{Lemme 2.4.3.} Pour tout schéma propre \(X\) sur un corps algébriquement clos \(k\) le groupe \(\Pic^0(X)\) est divisible par tout entier \(n\) premier à la caractéristique de \(k\).

\emph{Preuve :} Comme \(k\) est algébriquement clos, on a \(\Pic^0(X)\cong \Pic^0_{X/k}(k)\). Puisque \(\Pic^0_{X/k}\) est connexe et de type fini, l'endomorphisme puissance \(n\)-ème est surjectif lorsque \(n\) est premier à la caractéristique de \(k\) (SGA 3 XV 1.3 et VI 1.3.2). Mais \(k\) est algébriquement clos, donc l'endomorphisme correspondant sur le groupe \(\Pic^0_{X/k}(k)\) est surjectif, C.Q.F.D.

De 2.4.3 on conclut que dans la suite exacte
\[
0\longrightarrow\Pic^0(X)\longrightarrow\Pic^\tau(X)\longrightarrow T\longrightarrow0,
\]
\(\Pic^0(X)\) est divisible par un nombre premier \(p\) et \(T\) est un groupe de torsion, de sorte qu'on peut appliquer 2.1.3 à la restriction de l'application polynôme (2.3.1) à \(\Pic^\tau(X)\). On trouve ainsi :

